\section{Walk-Through: Local Labels Carry the Load}\label{sec:markov-walkthrough}

Section~\ref{sec:markov-interlude} described the scene-register
changepoint task and explained why a leaf summary that drops boundary
registers cannot recover the count of register shifts. We now ask the
empirical question: when local supervision is available (a leaf-level
label of the count, or of the boundary registers, on a sub-span), can
it stand in for some of the root labels? The models are a flat Fourier neural operator applied to
the full document and a tree of FNO leaves merged by a small learned
MLP, both trained against the same closed-form oracle on the synthetic
DGP of Section~\ref{sec:scene-register-example}; full architecture,
training, and DGP detail are in Appendix~\ref{app:mechanism-checks}.

The answer is yes, and Figure~\ref{fig:leaf-size-structural} carries
the headline. Each row fixes a root-label budget ($100\%$ down to
$10\%$); each column fixes a leaf size ($128$ tokens, where the tree
is a single leaf and degenerates to the flat baseline, down to $8$
tokens). Cell color is the tree-oracle root MAE on the held-out test
set, with darker cells meaning lower MAE. Reading right to left
across any row tells the central story: the single-leaf flat baseline
sits at root MAE $0.30$--$0.66$ across root shares, while the same
network applied as a tree of $32$- or $64$-token leaves drops below
$0.05$ MAE. The tree wins by an order of magnitude in the harder of
the two regimes we ran (twelve hidden registers, ten-to-twelve
segments per document). Reading down the figure traces what happens
when root labels become scarce: keeping the leaf size in the safe
band (32 to 64 tokens), the tree's MAE stays under $0.05$ all the way
down to a $20\%$ root-label budget, with most of the substituted
budget going to local labels at internal nodes. The failure regime
sits at the bottom-left of the same panel: at root share $10\%$ with
$8$-token leaves, MAE blows up to $0.66$, because the leaves are now
too small to carry the boundary information the audit requires. The
figure is both the positive result (label budget can move down the
tree) and its limit (leaves cannot get arbitrarily small).

The mechanism is the same boundary-state argument from
Section~\ref{sec:markov-interlude}: a leaf summary that knows its
first and last register can carry that information through a merge
that sums internal counts plus an indicator of register mismatch.
Local labels supply exactly the supervision the merge needs to learn
that boundary correction. Section~\ref{sec:hll-parity} verifies the
deterministic limit of the same construction (HyperLogLog recovered
byte-for-byte under the local laws), and Section~\ref{sec:manifesto-llm}
shows that the same down-the-tree label-budget logic survives when the
leaves are high-dimensional LLM summaries and the oracle is a
political-science rubric. Appendix~\ref{app:mechanism-checks} has the
DGP specification, the three supervision-allocation policies
(leaf-only, depth-equal, balanced-node), the recoverable-regime
companion plot, the per-cell budget-split heatmap, and the
architecture and training-stage details.

\begin{figure}[t]
    \centering
    \includegraphics[width=0.95\linewidth]{markov_leaf_size_fixed_structural.pdf}
    \caption{\textbf{Local labels carry the load: tree-oracle root MAE on the structural-regime synthetic DGP (twelve hidden registers, ten-to-twelve segments per document).} Each cell fixes a root-label budget (rows, $100\%$ to $10\%$) and a leaf size (columns, $128$ tokens down to $8$). The leaf-$128$ column is the single-leaf case where the tree degenerates to the flat FNO baseline; smaller leaves are the tree with local supervision at internal nodes. The safe band (leaves of $32$--$64$ tokens) holds MAE below $0.05$ across most root-label budgets. The failure regime at very small leaves and very low root share is visible in the bottom-left and is honest about where the framework stops working. The recoverable-regime companion is in Appendix~\ref{app:mechanism-checks}.}
    \label{fig:leaf-size-structural}
\end{figure}
